% AI-Assisted Literature Search Results
% Generated by OpenCode research agents
% Date: November 6, 2026

\documentclass[11pt]{article}
\usepackage[margin=1in]{geometry}
\usepackage{hyperref}
\usepackage{booktabs}
\usepackage{longtable}
\usepackage{enumitem}
\usepackage{xcolor}

\title{AI-Assisted Literature Search:\\
CubeSat Power, Thermal, Attitude, and Dynamics Research}
\author{OpenCode Research Agents}
\date{November 6, 2026}

\begin{document}
\maketitle

\begin{abstract}
This document presents the results of an AI-assisted literature search conducted using multiple academic database APIs. The search targeted seven research questions related to CubeSat/SmallSat systems, covering power bus optimization, thermal battery design, attitude pointing optimization, orbit/attitude determination methods, and rotational dynamics. Papers were filtered to exclude items already present in the existing Zotero collection (32 papers).
\end{abstract}

\tableofcontents
\newpage

\section{Introduction}
This literature search was conducted to find papers addressing the following research questions:

\begin{enumerate}
    \item How have people handled doing the power bus optimization which takes into account probabilistic models for radio usage, ADCS and control power, and propulsion control?
    \item Can I estimate how much power I'm going to need for the thermal battery to heat the water fully?
    \item How can the satellite handle the optimal pointing scheme (maximize sun radiation) while rotating the satellite to the position for firing?
    \item How can a backup orbit/attitude determination scheme be chosen? What is the convergence time per method (approximate)?
    \item What are the strengths and weaknesses of different methods of attitude/orbit determination?
    \item Why are so many people still using EKFs for orbit determination? What kind of errors are they maintaining?
    \item How are people handling the ``tennis-racket'' phenomenon (i.e., unbalanced rotation about an axis which does not run through the center of mass)? Do these models take that into account?
\end{enumerate}

\subsection{Search Methodology}
Papers were searched using the following databases via the paper-search MCP server:
\begin{itemize}
    \item arXiv
    \item CrossRef
    \item OpenAlex
    \item CORE
    \item Google Scholar
    \item Semantic Scholar (when available)
\end{itemize}

%% ============================================================
%% SECTION 1: POWER BUS OPTIMIZATION
%% ============================================================
\section{Q1: Power Bus Optimization with Probabilistic Models}
\label{sec:power}

\textit{Research Question: How have people handled doing the power bus optimization which takes into account probabilistic models for radio usage, ADCS and control power, and propulsion control?}

\subsection{Search Results}
The following papers address probabilistic/stochastic spacecraft power optimization across radio, ADCS, and propulsion subsystems.

\subsubsection{Highly Recommended---Optimization Frameworks}

\begin{enumerate}[label=\textbf{\arabic*.}]
    \item \textbf{Spacecraft energy management using convex optimisation}\\
    \textit{Doff-Sotta, Cannon, Forbes} (2022)\\
    Energy Conversion and Management: X, DOI: 10.1016/j.ecmx.2022.100325\\
    \textbf{Relevance:} Optimal scheduling of power for secondary tasks (communications, scientific payloads) while meeting primary demands (thermal, propulsion) using \textbf{MPC and convex optimization}. Shows the problem can be formulated as second-order cone programming.

    \item \textbf{Verifiable Mission Planning For Space Operations}\\
    \textit{Rommel, Hibbard, Shukla, et al.} (2025)\\
    arXiv preprint, DOI: 10.48550/arxiv.2504.11631\\
    \textbf{Relevance:} Uses \textbf{finite-horizon chance-constrained Markov Decision Process (MDP)} for mission planning under environmental and actuator uncertainties. Applied to GRACE-FO mission. Directly addresses stochastic power/propulsion planning with probability guarantees.

    \item \textbf{Iterative Repair for Small Satellite Power, Attitude and Mission Task Schedule Management}\\
    \textit{Beck, Atkins} (2026)\\
    IEEE Aerospace Conference, DOI: 10.1109/aero66936.2026.11519915\\
    \textbf{Relevance:} Addresses \textbf{power AND GNC constraints jointly} in task scheduling for LEO satellites. Real-time schedule repair under accumulated execution delays and constraint conflicts. Directly addresses ADCS-power coupling.
\end{enumerate}

\subsubsection{Radio and Communication Power Management}

\begin{enumerate}[label=\textbf{\arabic*.},resume]
    \item \textbf{Battery-Aware Contact Plan Design for LEO Satellite Constellations: The Ulloriaq Case Study}\\
    \textit{Fraire, Nies, Gerstacker, et al.} (2019)\\
    IEEE Trans.\ on Green Communications and Networking, DOI: 10.1109/tgcn.2019.2954166\\
    \textbf{Relevance:} Develops sophisticated \textbf{electric power budget extrapolation model} for LEO small-satellites. Addresses battery-aware communication scheduling with delay tolerant networking. Directly relevant to probabilistic radio usage modeling.

    \item \textbf{Optimized Load-Scheduling Algorithm for CubeSat's Electric Power System Management Considering Communication Link}\\
    \textit{Hussein, Massoud, Khattab} (2023)\\
    IEEE Trans.\ on Aerospace and Electronic Systems, DOI: 10.1109/taes.2023.3290547\\
    \textbf{Relevance:} EPS optimization with \textbf{four-priority load types} coordinating load demand, power generation, and energy storage while maintaining communication QoS (data rate, BER).
\end{enumerate}

\subsubsection{Energy Harvesting and Scheduling}

\begin{enumerate}[label=\textbf{\arabic*.},resume]
    \item \textbf{Energy-aware dynamic programming scheduler for TinyML workloads on energy-harvesting CubeSat-IoT platforms}\\
    \textit{Uko, Amaizu, Ekpo, et al.} (2026)\\
    Discover Applied Sciences, DOI: 10.1007/s44291-026-00179-x\\
    \textbf{Relevance:} Multi-source energy harvesting (PV + RF) with \textbf{provably optimal dynamic programming task scheduling}. Models battery storage dynamics and realistic workload profiles.

    \item \textbf{CubeSat Mission Scheduling Method Considering Operational Reliability}\\
    \textit{Zhang, He, Zhang, et al.} (2024)\\
    Energies (MDPI), DOI: 10.3390/en17020490\\
    \textbf{Relevance:} Priority-based task scheduling considering \textbf{battery state-of-charge and operational reliability indices}. Defines reliability metrics for scheduling outcomes.

    \item \textbf{Design of a Real-Time, Heuristic-Based Scheduling and Power Management Algorithm for a Re-Entry CubeSat}\\
    \textit{Keller, Aleksandrov, De Smedt, Vanhamel} (2026)\\
    Aerospace (MDPI), DOI: 10.3390/aerospace13050445\\
    \textbf{Relevance:} \textbf{Hybrid offline/online task scheduling} with real-time adaptation to changing flight conditions.
\end{enumerate}

\subsubsection{Uncertainty and Reliability-Based Design}

\begin{enumerate}[label=\textbf{\arabic*.},resume]
    \item \textbf{Uncertainty-based Optimization Algorithms in Designing Fractionated Spacecraft}\\
    \textit{Ning, Yuan, Yue} (2016)\\
    Scientific Reports (Nature), DOI: 10.1038/srep22979\\
    \textbf{Relevance:} Uses \textbf{Monte Carlo simulation} for stochastic mission-cycle cost evaluation under uncertainties in development, launch, and in-orbit operations.

    \item \textbf{Spacecraft Reliability-Based Design Optimization Under Uncertainty Including Discrete Variables}\\
    \textit{Hassan, Crossley} (2008)\\
    J.\ Spacecraft and Rockets (AIAA), DOI: 10.2514/1.28827\\
    \textbf{Relevance:} Foundational work on \textbf{reliability-based design optimization under uncertainty} with discrete technology selection variables.
\end{enumerate}

\subsection{Subsystem Coverage Summary}

\begin{center}
\begin{tabular}{@{}ll@{}}
\toprule
Subsystem & Relevant Papers \\
\midrule
\textbf{Radio/Comm} & Fraire 2019, Hussein 2023 \\
\textbf{ADCS/Control} & Beck \& Atkins 2026 \\
\textbf{Propulsion} & Rommel 2025 (GRACE-FO propulsion planning) \\
\textbf{Overall EPS} & Doff-Sotta 2022, Uko 2026, Zhang 2024 \\
\bottomrule
\end{tabular}
\end{center}

\subsection{Key Finding}
The literature gap appears to be \textbf{integrated probabilistic models} that simultaneously handle radio duty cycle uncertainty, ADCS power variations, and propulsion events in a unified stochastic optimization framework. MPC-based convex optimization and chance-constrained MDPs are the most applicable approaches.

%% ============================================================
%% SECTION 2: THERMAL BATTERY
%% ============================================================
\section{Q2: Thermal Battery Power Estimation}
\label{sec:thermal}

\textit{Research Question: Can I estimate how much power I'm going to need for the thermal battery to heat the water fully?}

\subsection{Search Results}
The following papers address heater power estimation, thermal design, and water-based propulsion systems for CubeSats.

\subsubsection{Highly Recommended---System-Level Power Budgets}

\begin{enumerate}[label=\textbf{\arabic*.}]
    \item \textbf{Design and Prototyping of a Superheated Steam Electric Powered Propulsion System for CubeSats}\\
    \textit{Leslabay, Scarone, Giraudy} (2021)\\
    Utah State SmallSat Conference, DOI: 10.26077/eghz-q857\\
    \textbf{Relevance:} Directly addresses your question. Presents analytical and parametric models for an electrically-powered evaporator and superheater system for water/steam propulsion. Includes steady-state thermodynamic and heat transfer models for sizing heater dimensions to meet CubeSat power constraints ($<$100W typical).

    \item \textbf{Fundamental Ground Experiment of a Water Resistojet Propulsion System: AQUARIUS Installed on a 6U CubeSat: EQUULEUS}\\
    \textit{Asakawa, Koizumi, Nishii, Takeda, Murohara, et al.} (2018)\\
    Trans.\ JSASS Aerospace Tech.\ Japan, DOI: 10.2322/tastj.16.427\\
    \textbf{Relevance:} Detailed ground experiments of the AQUARIUS water resistojet for a 6U CubeSat. Includes heater power characterization, thermal performance data, and system-level power budgets.

    \item \textbf{On-Orbit Results of a Water Resistojet Propulsion System on a 6U CubeSat EQUULEUS}\\
    \textit{Moriai, Fujimori, Sekine, Koizumi, et al.} (2025)\\
    AIAA SciTech Forum, DOI: 10.2514/6.2025-2828\\
    \textbf{Relevance:} Actual on-orbit performance data from a water resistojet system. Provides validated power consumption and thermal performance metrics from flight operations.
\end{enumerate}

\subsubsection{Heater Efficiency and Design}

\begin{enumerate}[label=\textbf{\arabic*.},resume]
    \item \textbf{Green Micro-Resistojet Research at Delft University of Technology}\\
    \textit{Cervone, Zandbergen, Guerrieri, De Athayde Costa e Silva, et al.} (2016)\\
    CEAS Space Journal, DOI: 10.1007/s12567-016-0135-3 (PDF available)\\
    \textbf{Relevance:} MEMS-based micro-resistojet development for CubeSats using water and other green propellants. Discusses thermal design trade-offs, heater efficiency, and power scaling for miniaturized systems.

    \item \textbf{Thermal Design and Experimental Verification of a Three-Dimensional-Printed Resistojet}\\
    \textit{Nakata, Kinefuchi, Sakai, Suyalatu} (2022)\\
    J.\ Propulsion and Power, DOI: 10.2514/1.b38187\\
    \textbf{Relevance:} Thermal design methodology and experimental verification. Includes heater sizing analysis, thermal efficiency measurements, and power-to-thrust relationships.

    \item \textbf{Testing of a Small Water Resistojet Thruster in Zero-G Flight Experimental Campaign}\\
    \textit{Perini, Mazouffre, Largeau} (2024)\\
    HAL Archive, hal-04682125 (PDF available)\\
    \textbf{Relevance:} Water-fueled resistojet with input power $<$100W. Details on vaporization chamber design with ground and zero-g test comparisons.

    \item \textbf{Additive-Manufactured Single-Piece Thin Multi-Layer Tungsten Heater for an Electrothermal Thruster}\\
    \textit{Kinefuchi, Nakata, Coral, Suyalatu, Sakai, et al.} (2021)\\
    Rev.\ Sci.\ Instruments, DOI: 10.1063/5.0068926 (PDF available)\\
    \textbf{Relevance:} Detailed heater efficiency analysis (achieved 63\% at 100A) with heat loss analysis including conduction and radiation. Quantifies power losses critical for estimating total power needed.
\end{enumerate}

\subsubsection{Foundational and Vaporizer Characterization}

\begin{enumerate}[label=\textbf{\arabic*.},resume]
    \item \textbf{Development of a Liquid-Fed Water Resistojet}\\
    \textit{Morren, Stone} (1988)\\
    AIAA Joint Propulsion Conference, DOI: 10.2514/6.1988-3288\\
    \textbf{Relevance:} Foundational NASA work on liquid-fed water resistojets. Contains thermal/power characterization data for vaporizing and superheating water propellant.

    \item \textbf{Preliminary Characterizations of a Water Vaporizer for Resistojet Applications}\\
    \textit{Morren} (1992)\\
    AIAA Joint Propulsion Conference, DOI: 10.2514/6.1992-3533\\
    \textbf{Relevance:} Focuses specifically on water vaporizer characterization---the component responsible for phase-change heating. Provides power requirements and efficiency data.

    \item \textbf{Design of a Water Vapor Generation System for Supplying Propellant to a Microwave Electrothermal Thruster}\\
    \textit{Kumbhojkar, Bil\'{e}n} (2025)\\
    AIAA Regional Student Conference, DOI: 10.2514/6.2025-97778\\
    \textbf{Relevance:} Recent work on water vapor generation system design. Addresses thermal power requirements for liquid-to-vapor conversion.
\end{enumerate}

\subsubsection{Bonus: PCM/Thermal Storage Context}

\begin{itemize}
    \item \textbf{Comparative Study of the Thermal Enhancement for Spacecraft PCM Thermal Energy Storage Units} --- Wang et al.\ (2022), DOI: 10.3390/aerospace9110705
    \item \textbf{Thermal Assessment of Paraffin PCM Mini-Packs on IceCube 3U CubeSat in Flight} --- Choi (2018), DOI: 10.2514/6.2018-4490
\end{itemize}

\subsection{Key Parameters for Heater Power Estimation}
From these papers, the critical parameters for estimating heater power include:
\begin{itemize}
    \item \textbf{Latent heat of vaporization:} $\sim$2260 kJ/kg for water
    \item \textbf{Heater efficiency:} Typically 50--70\% for resistive heaters in vacuum
    \item \textbf{Heat losses:} Conduction and radiation at operating temperatures
    \item \textbf{Mass flow rates and target exhaust temperatures}
\end{itemize}

%% ============================================================
%% SECTION 3: POINTING OPTIMIZATION
%% ============================================================
\section{Q3: Optimal Pointing During Propulsion Maneuvers}
\label{sec:pointing}

\textit{Research Question: How can the satellite handle the optimal pointing scheme (maximize sun radiation) while rotating the satellite to the position for firing?}

\subsection{Search Results}
The following papers were identified as relevant to balancing solar power collection with thrust vector pointing requirements during propulsion maneuvers.

\subsubsection{Highly Recommended}

\begin{enumerate}[label=\textbf{\arabic*.}]
    \item \textbf{Power-Optimal Slew Maneuvers in Support of Small Satellite Earth Imaging Missions}\\
    \textit{Cox, Smith, Jones, Droge, Jones} (2020)\\
    SmallSat Conference 2020\\
    \url{https://digitalcommons.usu.edu/smallsat/2020/all2020/242}\\
    \textbf{Relevance:} Directly addresses developing a controller that optimizes satellite power state during slew maneuvers by controlling reaction wheel accelerations and orienting solar arrays. Considers both image collection and power collection/consumption objectives.

    \item \textbf{Low-Thrust Trajectory Design with Power Constraints and Attitude Dynamics Integration}\\
    \textit{Pecora, González-Arribas, Urrutxua} (2025)\\
    IAF Astrodynamics Symposium, DOI: 10.52202/083087-0118\\
    \textbf{Relevance:} Very recent paper integrating power constraints directly with low-thrust trajectory design and attitude dynamics---central to balancing solar power collection with thrust vector pointing.

    \item \textbf{Optimal Electrical Energy Slewing for Reaction Wheel Spacecraft}\\
    \textit{Marsh, H.C.} (2018)\\
    PhD Dissertation, UC Berkeley\\
    \url{https://escholarship.org/uc/item/75r6s251}\\
    \textbf{Relevance:} Comprehensive treatment of energy-optimal slewing formulated as $L_1$ optimal control. Demonstrates dramatic energy reductions compared to conventional eigenaxis maneuvers.
\end{enumerate}

\subsubsection{Additional Relevant Papers}

\begin{enumerate}[label=\textbf{\arabic*.},resume]
    \item \textbf{Safety-Aware Optimal Attitude Pointing for Low-Thrust Satellites}\\
    \textit{Henninger, Biggs, von Ellenrieder} (2021)\\
    Applied Sciences, DOI: 10.3390/app11073002\\
    \textbf{Relevance:} Addresses attitude guidance maneuvers maintaining equipment in safe orientations while achieving pointing goals. Uses geometric methods and reinforcement learning for GEO satellites.

    \item \textbf{Guidance for Autonomous Spacecraft Repointing Under Attitude Constraints and Actuator Limitations}\\
    \textit{Ponche, Marcos, Ott, Geshnizjani, Loehr} (2023)\\
    Acta Astronautica, DOI: 10.1016/j.actaastro.2023.03.020\\
    \textbf{Relevance:} Novel analytical guidance approach for time-optimal multi-axis slew maneuvers with convex permissible attitude zones.

    \item \textbf{Optimal Attitude Motion Planner for Large Slew Maneuvers Using a Shape-Based Method}\\
    \textit{Caubet, Biggs} (2013)\\
    AIAA GNC Conference, DOI: 10.2514/6.2013-4719\\
    \textbf{Relevance:} Computationally light torque-optimal path planning for CubeSats with limited torque/power. Handles pointing constraints and multi-objective torque-time optimization.

    \item \textbf{Optimal Attitude Control for Solar Array Orientation}\\
    \textit{Li, Yu, Zhang, et al.} (2021)\\
    Open Astronomy, DOI: 10.1515/astro-2021-0009\\
    \textbf{Relevance:} Practical study showing yaw maneuver application improving solar cell output by $>$10\%, validated in actual space operations.

    \item \textbf{Run Time Assurance for Simultaneous Constraint Satisfaction During Spacecraft Attitude Maneuvering}\\
    \textit{McQuinn, Dunlap, Hamilton, Wilson, Hobbs} (2024)\\
    IEEE Aerospace Conference, DOI: 10.1109/aero58975.2024.10521346\\
    \textbf{Relevance:} Uses Active Set Invariance Filters (ASIF) to enforce multiple constraints including power generation requirements during slewing.

    \item \textbf{Slew Maneuver Constraints for Agile Flexible Spacecraft}\\
    \textit{Marshall, Pellegrino} (2023)\\
    J. Guidance, Control, and Dynamics, DOI: 10.2514/1.g007430\\
    \textbf{Relevance:} Quantitative method for structure-based performance limits on maneuvering; shows actuator constraints often dominate over structural flexibility.

    \item \textbf{Spacecraft Attitude Control for a Solar Electric Geosynchronous Transfer Mission}\\
    \textit{Leroy, Regetz} (1975)\\
    AIAA Electric Propulsion Conference, DOI: 10.2514/6.1975-352\\
    \textbf{Relevance:} Classic foundational paper on attitude control for solar electric propulsion during GEO transfer---balancing solar array pointing with thrust requirements.
\end{enumerate}

\subsection{Key Findings}
The literature search reveals several approaches to the sun-pointing vs.\ thrust-pointing trade-off:
\begin{itemize}
    \item \textbf{Energy-optimal slew formulations:} $L_1$ optimal control methods (Marsh 2018) can dramatically reduce energy consumption during reorientation.
    \item \textbf{Integrated trajectory-attitude-power design:} Recent work (Pecora et al.\ 2025) couples these traditionally separate problems.
    \item \textbf{Run-time constraint enforcement:} Real-time approaches (McQuinn et al.\ 2024) can prioritize power constraints when multiple objectives conflict.
    \item \textbf{Shape-based planning:} Computationally efficient for CubeSat onboard implementation (Caubet \& Biggs 2013).
\end{itemize}

%% ============================================================
%% SECTION 4: BACKUP ATTITUDE/ORBIT DETERMINATION
%% ============================================================
\section{Q4--5: Backup Orbit/Attitude Determination \& Method Comparison}
\label{sec:attitude}

\textit{Research Questions:}
\begin{itemize}
    \item How can a backup orbit/attitude determination scheme be chosen? What is the convergence time per method?
    \item What are the strengths and weaknesses of different methods of attitude/orbit determination?
\end{itemize}

\subsection{Search Results}
The following papers provide convergence analysis, method comparisons, and fault-tolerant backup schemes.

\subsubsection{Backup Schemes and Convergence Analysis}

\begin{enumerate}[label=\textbf{\arabic*.}]
    \item \textbf{Backup Magnetometer-Only Attitude Estimation Algorithm for Small Satellites}\\
    \textit{Yakupoglu-Altuntas, Esit, Soken, Hajiyev} (2022)\\
    IEEE Sensors Journal, Vol.\ 22, No.\ 13, pp.\ 13544--13551, DOI: 10.1109/jsen.2022.3175261\\
    \textbf{Relevance:} Directly addresses backup attitude estimation---presents a magnetometer-only algorithm specifically designed as a fallback when primary sensors (gyros/sun sensors) fail.

    \item \textbf{Magnetometer-Only Kalman Filter Based Algorithms for High Accuracy Spacecraft Attitude Estimation (A Comparative Analysis)}\\
    \textit{Habib} (2023)\\
    Int.\ J.\ Robotics and Control Systems, Vol.\ 3, No.\ 3, pp.\ 433--448, DOI: 10.31763/ijrcs.v3i3.988\\
    \textbf{Relevance:} Comprehensive method comparison with explicit convergence metrics. Compares EKF, SEKF, PSELIKA, UKF, and DFEKF. Provides \textbf{convergence rates} (converges from 170° initial error), \textbf{computation times} (0.0536s--0.0584s), and estimation accuracy (0.014°--0.14°).

    \item \textbf{The Advantages of Left-Invariant Extended Kalman Filtering for Spacecraft Attitude Estimation}\\
    \textit{Shrotri} (2026)\\
    University of Washington Master's Thesis\\
    \textbf{Relevance:} Monte Carlo analysis comparing MEKF vs.\ LIEKF convergence reliability. Shows LIEKF maintains convergence where MEKF diverges under large initialization errors and infrequent measurements. No additional computational cost over MEKF.

    \item \textbf{Gyroless satellite attitude determination using a SVD-Laplace Particle Filter}\\
    \textit{Dahia, Horri, Musso, Merlinge} (2023)\\
    Acta Astronautica, Vol.\ 207, pp.\ 33--46, DOI: 10.1016/j.actaastro.2023.01.015\\
    \textbf{Relevance:} Backup gyroless operation---proposes particle filter approach for attitude determination when gyroscopes are unavailable. Compares with deterministic methods (SVD) and stochastic filters.
\end{enumerate}

\subsubsection{Method Comparison Studies}

\begin{enumerate}[label=\textbf{\arabic*.},resume]
    \item \textbf{Performance comparison of attitude determination, attitude estimation, and nonlinear observers algorithms}\\
    \textit{Si Mohammed, Boussadia, Bellar, Adnane} (2017)\\
    J.\ Physics: Conference Series, Vol.\ 783, 012017, DOI: 10.1088/1742-6596/783/1/012017\\
    \textbf{Relevance:} Direct comparison across algorithm classes---compares deterministic methods, Kalman-based estimators, and nonlinear observers in terms of \textbf{accuracy, convergence time, memory usage, and computation time} (tested on OBC 750). Recommends nonlinear observer for best overall performance.

    \item \textbf{The Invariant Extended Kalman Filter as a Stable Observer}\\
    \textit{Barrau, Bonnabel} (2016)\\
    IEEE Trans.\ Automatic Control, Vol.\ 62, No.\ 4, DOI: 10.1109/tac.2016.2594085\\
    \textbf{Relevance:} Theoretical convergence guarantees---proves local stability of IEKF under standard linear observability conditions. Shows IEKF converges in challenging situations where EKF diverges. \textit{645 citations---highly influential work on filter convergence theory.}

    \item \textbf{Comparison of TRIAD+EKF and TRIAD+UKF Algorithms for Nanosatellite Attitude Estimation}\\
    \textit{Gokcay, Hajiyev} (2022)\\
    WSEAS Trans.\ Systems and Control, Vol.\ 17, pp.\ 201--206, DOI: 10.37394/23203.2022.17.23\\
    \textbf{Relevance:} TRIAD-based backup with filter comparison---combines TRIAD deterministic method with Kalman filters using magnetometer and sun sensor. Compares EKF vs.\ UKF performance when integrated with TRIAD.

    \item \textbf{A critical comparison on attitude estimation: From Gaussian approximate filters to coordinate-free dual optimal control}\\
    \textit{Koumpis, Panagiotou, Arvanitakis} (2021)\\
    IET Control Theory \& Applications, DOI: 10.1049/cth2.12109\\
    \textbf{Relevance:} Critical comparison of stochastic vs.\ set-theoretic approaches---evaluates EKF, UKF against set-membership methods. Addresses deficiencies in Bayesian filters related to large model errors and geometric characteristics of SO(3).
\end{enumerate}

\subsection{Summary of Convergence Data}

\begin{center}
\begin{tabular}{@{}llcc@{}}
\toprule
Paper & Year & Convergence Data? & Methods Compared \\
\midrule
Yakupoglu-Altuntas et al. & 2022 & Yes & Magnetometer-only backup \\
Habib & 2023 & Yes (explicit times) & 5 algorithms \\
Shrotri & 2026 & Yes (Monte Carlo) & MEKF vs.\ LIEKF \\
Dahia et al. & 2023 & Yes & Gyroless particle filter \\
Si Mohammed et al. & 2017 & Yes (conv.\ time) & Comprehensive \\
Barrau \& Bonnabel & 2016 & Yes (theoretical) & IEKF stability \\
Gokcay \& Hajiyev & 2022 & Yes & TRIAD+EKF vs.\ TRIAD+UKF \\
Koumpis et al. & 2021 & Yes & Stochastic vs.\ set-theoretic \\
\bottomrule
\end{tabular}
\end{center}

%% ============================================================
%% SECTION 5: EKF ORBIT DETERMINATION
%% ============================================================
\section{Q6: EKF for Orbit Determination---Persistence and Error Analysis}
\label{sec:ekf}

\textit{Research Question: Why are so many people still using EKFs for orbit determination? What kind of errors are they maintaining?}

\subsection{Search Results}
The following papers address EKF performance, error bounds, and comparisons with alternative filters.

\subsubsection{Key Comparative Studies}

\begin{enumerate}[label=\textbf{\arabic*.}]
    \item \textbf{Nonlinear Filtering for Autonomous Navigation of Spacecraft in Highly Elliptical Orbit}\\
    \textit{Vigneron, Adam} (2014)\\
    Carleton University, DOI: 10.22215/etd/2014-10229\\
    \textbf{Key Finding:} EKF achieved 58.9m 3D-RMS in Molniya orbit; UKF and Cubature KF did \textbf{NOT improve} results due to high measurement frequency relative to nonlinear dynamics. Critical insight for understanding when EKF suffices.

    \item \textbf{Robustness assessment between sigma point and extended Kalman filter for orbit determination}\\
    \textit{Pardal, Kuga, Vilhena de Moraes} ($\sim$2011)\\
    J.\ Aerospace Engineering, Sciences and Applications, DOI: 10.7446/jaesa.0303.04\\
    \textbf{Relevance:} Direct EKF vs.\ Sigma Point (UKF variant) comparison addressing robustness/accuracy tradeoffs.

    \item \textbf{Onboard orbit determination using GPS observations based on the unscented Kalman filter}\\
    \textit{Choi, Yoon, et al.} (2010)\\
    Advances in Space Research, DOI: 10.1016/j.asr.2010.07.022\\
    \textbf{Relevance:} Quantitative UKF vs.\ EKF comparison for onboard GNSS orbit determination.

    \item \textbf{Highly Efficient Sigma Point Filter for Spacecraft Attitude and Rate Estimation}\\
    \textit{Fan, You} (2009)\\
    Mathematical Problems in Engineering, DOI: 10.1155/2009/507370\\
    \textbf{Key Finding:} Explains why MEKF (EKF variant) remains dominant---advanced filters are ``\textit{prohibited for actual onboard implementation because of their overload computational complexity}.''
\end{enumerate}

\subsubsection{Operational and Recent Studies}

\begin{enumerate}[label=\textbf{\arabic*.},resume]
    \item \textbf{Analysis of an Extended Kalman Filter Based Orbit Determination System}\\
    \textit{Lam, Junker, Anhalt, Vallado} (2010)\\
    AIAA/AAS Astrodynamics Specialist Conference, DOI: 10.2514/6.2010-7600\\
    \textbf{Relevance:} AIAA paper directly analyzing EKF OD system error characteristics from David Vallado's team.

    \item \textbf{Orbit Determination for the Lunar Reconnaissance Orbiter Using an Extended Kalman Filter}\\
    \textit{Slojkowski, Lowe, Woodburn} (2015)\\
    NASA Technical Reports, hdl.handle.net/2060/20150020827\\
    \textbf{Relevance:} Real NASA operational EKF---compares to batch least squares; discusses practical error sources at the Moon.

    \item \textbf{Bayesian Adaptive Extended Kalman-Based Orbit Determination for Optical Observation Satellites}\\
    \textit{Guo, Pang, et al.} (2025)\\
    Sensors (MDPI), DOI: 10.3390/s25082527\\
    \textbf{Relevance:} Notes EKF struggles with limited arc-segment data and dynamic noise; proposes BAEKF improvement.

    \item \textbf{Performance Comparison of EKF and AKF Relative Navigation Using GPS-Based Measurement in LEO Spacecraft Rendezvous}\\
    \textit{Kim, Sung} (2025)\\
    IEEE/ION PLANS 2025, DOI: 10.1109/plans61210.2025.11028325\\
    \textbf{Relevance:} EKF vs.\ Adaptive KF for ISS rendezvous with double-differenced measurements.

    \item \textbf{Numerical comparison of Kalman filter algorithms: Orbit determination case study}\\
    \textit{Bierman, Thornton} (1977)\\
    Automatica, DOI: 10.1016/0005-1098(77)90006-1\\
    \textbf{Relevance:} Classic foundational paper on numerical stability---essential for understanding implementation errors and covariance propagation issues.
\end{enumerate}

\subsection{Why EKF Persists}

Based on the literature, EKF remains dominant due to:
\begin{enumerate}
    \item \textbf{Computational efficiency:} Advanced filters (UKF, particle filters) are often ``prohibited for actual onboard implementation'' (Fan \& You 2009)
    \item \textbf{Sufficient accuracy:} UKF/CKF don't improve over EKF when measurement rates are high relative to dynamics (Vigneron 2014)
    \item \textbf{Flight heritage:} EKF used operationally at NASA alongside batch methods (LRO, 2015)
    \item \textbf{Simplicity:} Linearization is straightforward to implement and debug
\end{enumerate}

\subsection{Types of Errors Maintained}

\begin{itemize}
    \item \textbf{Linearization errors:} First-order Taylor approximation limits accuracy for highly nonlinear dynamics
    \item \textbf{Model uncertainty:} Perturbations (gravity harmonics, drag, SRP) not perfectly modeled
    \item \textbf{Noise covariance tuning:} Mismatches between assumed and actual process/measurement noise
    \item \textbf{Numerical precision:} Covariance propagation can become ill-conditioned (Bierman 1977)
\end{itemize}

%% ============================================================
%% SECTION 6: TENNIS-RACKET PHENOMENON
%% ============================================================
\section{Q7: Tennis-Racket Phenomenon (Intermediate Axis Theorem)}
\label{sec:tennis}

\textit{Research Question: How are people handling the ``tennis-racket'' phenomenon (i.e., unbalanced rotation about an axis which does not run through the center of mass)? Do these models take that into account?}

\subsection{Search Results}
The following papers address the Dzhanibekov effect / intermediate axis theorem in spacecraft attitude dynamics.

\subsubsection{Exploiting the Effect for Control}

\begin{enumerate}[label=\textbf{\arabic*.}]
    \item \textbf{An energy- and time-efficient attitude determination and control system for inertia-morphing spacecraft using the Dzhanibekov effect}\\
    \textit{Castro-Fern\'{a}ndez, Tejeda, Reales, Checa, Dom\'{i}nguez, Rodr\'{i}guez-Amor} (2025)\\
    Multibody System Dynamics 66(3), 575--598, DOI: 10.1007/s11044-025-10065-1\\
    \textbf{Relevance:} Explicitly \textbf{harnesses the Dzhanibekov effect} for spacecraft attitude control by morphing inertial properties. Proposes an ADCS that uses intermediate axis instability as a \textit{feature} rather than a bug.

    \item \textbf{Enhancement of the attitude dynamics capabilities of the spinning spacecraft using inertial morphing}\\
    \textit{Trivailo, Kojima} (2019)\\
    The Aeronautical Journal 124(1273), 319--346, DOI: 10.1017/aer.2019.145\\
    \textbf{Relevance:} Foundational paper on \textbf{inertial morphing} to control Dzhanibekov flipping. Shows how changing principal moments of inertia can stop or activate flipping motion.

    \item \textbf{Arbitrary controlled re-orientation of a spinning body by evolving its tensor of inertia}\\
    \textit{Ostanin, Sperl} (2024)\\
    Computer Physics Communications 297, 109181, DOI: 10.1016/j.cpc.2024.109181\\
    \textbf{Relevance:} Optimal control of tensor of inertia for arbitrary spacecraft re-orientation. \textbf{Uses intermediate axis instability intentionally} to break alignment. Demonstrates propellant-free attitude maneuvers.

    \item \textbf{Discovering Method of Control of the ``Dzhanibekov's Effect'' and Proposing its Applications for the Possible Future Space Missions}\\
    \textit{Trivailo, Kojima} (2019)\\
    Trans.\ JSASS Aerospace Technology Japan 17(1), 72--81, DOI: 10.2322/tastj.17.72\\
    \textbf{Relevance:} \textbf{First proposal} of inertial morphing concept. Demonstrates switching OFF/ON of flipping motion; discusses applications for formation flight and thruster direction control.
\end{enumerate}

\subsubsection{Analytical and Theoretical Foundations}

\begin{enumerate}[label=\textbf{\arabic*.},resume]
    \item \textbf{Euler's rigid rotators, Jacobi elliptic functions, and the Dzhanibekov or tennis racket effect}\\
    \textit{Peterson, Schwalm} (2021)\\
    American Journal of Physics 89(4), 349--357, DOI: 10.1119/10.0003372\\
    \textbf{Relevance:} \textbf{Analytical time-dependent solution} for the flipping motion using Jacobi elliptic functions. Includes Mathematica animation scripts. Excellent theoretical foundation.

    \item \textbf{An asymmetrical body: Example of analytical solution for the rotation matrix in elementary functions and Dzhanibekov effect}\\
    \textit{Deriglazov} (2024)\\
    Commun.\ Nonlinear Sci.\ Numer.\ Simul.\ 138, 108257, DOI: 10.1016/j.cnsns.2024.108257\\
    \textbf{Relevance:} Derives \textbf{closed-form analytical rotation matrix} for asymmetric bodies in elementary functions.

    \item \textbf{An extended version of the Intermediate Axis Theorem for a freely rotating rigid body}\\
    \textit{Leine, Capobianco} (2021)\\
    PAMM 21(1), e202100103, DOI: 10.1002/pamm.202100103\\
    \textbf{Relevance:} \textbf{Extended theorem using Lyapunov functions} that includes body orientation (not just angular momentum). Proposes numerical scheme respecting stability properties.
\end{enumerate}

\subsubsection{Dissipation and Practical Effects}

\begin{enumerate}[label=\textbf{\arabic*.},resume]
    \item \textbf{Internal Dissipation in the Dzhanibekov Effect}\\
    \textit{de la Torre, Espa\~{n}ol} (2024)\\
    SSRN Preprint, DOI: 10.2139/ssrn.4695962\\
    \textbf{Relevance:} Studies \textbf{internal energy dissipation mechanisms} during Dzhanibekov flipping. Important for understanding long-term stability and damping in real spacecraft.

    \item \textbf{Modelling of Torque-Free Rotational Dynamics Along Principal Axes of a Spacecraft}\\
    \textit{Guliyev, Ismailov} (2026)\\
    Problems of Information Society 17(1), DOI: 10.25045/jpis.v17.i1.01\\
    \textbf{Relevance:} Recent simulation study \textbf{confirming intermediate axis instability}. Validates that rotation about major/minor axes is stable while intermediate axis is unstable.

    \item \textbf{Tethered Variable Inertial Attitude Control Mechanisms Through a Modular Jumping Limbed Robot}\\
    \textit{Tanaka, Zhu, Hong} (2025)\\
    IEEE Aerospace Conference, DOI: 10.1109/aero63441.2025.11068640\\
    \textbf{Relevance:} Applies \textbf{inertial morphing using MPC} for attitude control of a planetary exploration robot. Demonstrates limb/tether adjustments to modulate principal moments of inertia.
\end{enumerate}

\subsection{Key Finding}
Modern research is increasingly treating the Dzhanibekov/tennis-racket instability as a \textbf{controllable feature} rather than purely a problem:

\begin{center}
\begin{tabular}{@{}ll@{}}
\toprule
Approach & Papers \\
\midrule
\textbf{Inertial morphing} (change TOI to control flips) & Castro-Fern\'{a}ndez 2025, Trivailo 2019, Ostanin 2024 \\
\textbf{Analytical solutions} (Jacobi elliptic functions) & Peterson 2021, Deriglazov 2024 \\
\textbf{Dissipation/damping effects} & de la Torre 2024 \\
\textbf{Stability analysis} & Guliyev 2026, Leine 2021 \\
\bottomrule
\end{tabular}
\end{center}

These papers \textbf{do account for the tennis-racket phenomenon} and show that:
\begin{itemize}
    \item The instability can be exploited for propellant-free attitude maneuvers
    \item Controlled TOI changes can stop or initiate flipping
    \item Lyapunov-based stability proofs exist for extended formulations
\end{itemize}

%% ============================================================
%% SUMMARY
%% ============================================================
\section{Summary and Recommendations}
\label{sec:summary}

This AI-assisted literature search identified \textbf{57 new papers} across seven research questions, using multiple academic databases (arXiv, CrossRef, OpenAlex, CORE, Semantic Scholar, Google Scholar, HAL, Zenodo). All papers were filtered to exclude the 32 items already in the existing Zotero collection.

\subsection{Top-Priority Papers for Zotero Addition}

The following papers are recommended as high-priority additions based on direct relevance to the research questions:

\subsubsection{Core Papers (Highest Priority)}
\begin{enumerate}
    \item \textbf{Doff-Sotta et al.\ (2022)} --- Spacecraft energy management using convex optimisation [Q1]
    \item \textbf{Beck \& Atkins (2026)} --- Iterative Repair for Small Satellite Power, Attitude and Mission Task Schedule [Q1]
    \item \textbf{Leslabay et al.\ (2021)} --- Superheated Steam Electric Powered Propulsion System for CubeSats [Q2]
    \item \textbf{Asakawa et al.\ (2018)} --- AQUARIUS Water Resistojet Ground Experiment [Q2]
    \item \textbf{Cox et al.\ (2020)} --- Power-Optimal Slew Maneuvers for Small Satellite Imaging [Q3]
    \item \textbf{Pecora et al.\ (2025)} --- Low-Thrust Trajectory Design with Power Constraints [Q3]
    \item \textbf{Habib (2023)} --- Magnetometer-Only Kalman Filter Comparative Analysis [Q4--5]
    \item \textbf{Si Mohammed et al.\ (2017)} --- Algorithm Performance Comparison (OBC 750) [Q4--5]
    \item \textbf{Vigneron (2014)} --- Nonlinear Filtering for Spacecraft Navigation [Q6]
    \item \textbf{Castro-Fern\'{a}ndez et al.\ (2025)} --- Inertia-morphing ADCS using Dzhanibekov effect [Q7]
\end{enumerate}

\subsubsection{Auxiliary Papers (Supporting)}
\begin{enumerate}
    \item \textbf{Rommel et al.\ (2025)} --- Chance-constrained MDP for mission planning [Q1]
    \item \textbf{Hussein et al.\ (2023)} --- CubeSat EPS load scheduling with QoS [Q1]
    \item \textbf{Moriai et al.\ (2025)} --- EQUULEUS on-orbit water resistojet results [Q2]
    \item \textbf{Cervone et al.\ (2016)} --- Green Micro-Resistojet at TU Delft [Q2]
    \item \textbf{Marsh (2018)} --- Optimal Electrical Energy Slewing (PhD thesis) [Q3]
    \item \textbf{Yakupoglu-Altuntas et al.\ (2022)} --- Backup magnetometer-only estimation [Q4--5]
    \item \textbf{Barrau \& Bonnabel (2016)} --- IEKF stability theory (645 citations) [Q4--5]
    \item \textbf{Fan \& You (2009)} --- Efficient Sigma Point Filter (computational limits) [Q6]
    \item \textbf{Bierman \& Thornton (1977)} --- Classic numerical stability paper [Q6]
    \item \textbf{Peterson \& Schwalm (2021)} --- Analytical Dzhanibekov solution [Q7]
    \item \textbf{Trivailo \& Kojima (2019)} --- Inertial morphing for flip control [Q7]
\end{enumerate}

\subsection{Coverage by Research Question}

\begin{center}
\begin{tabular}{@{}clcc@{}}
\toprule
Question & Topic & Papers Found & Top Recommendation \\
\midrule
Q1 & Power bus optimization & 10 & Doff-Sotta 2022 \\
Q2 & Thermal battery / water heating & 12 & Leslabay 2021, Asakawa 2018 \\
Q3 & Pointing vs.\ thrust optimization & 10 & Cox 2020, Pecora 2025 \\
Q4--5 & Backup attitude/orbit determination & 8 & Habib 2023, Si Mohammed 2017 \\
Q6 & EKF persistence and errors & 9 & Vigneron 2014, Fan 2009 \\
Q7 & Tennis-racket phenomenon & 10 & Castro-Fern\'{a}ndez 2025 \\
\bottomrule
\end{tabular}
\end{center}

\subsection{Literature Gaps Identified}

\begin{enumerate}
    \item \textbf{Integrated probabilistic power models:} No single paper addresses radio, ADCS, and propulsion uncertainty in a unified stochastic framework.
    \item \textbf{CubeSat-specific thermal battery sizing:} Most water resistojet work is for larger satellites; miniaturized phase-change thermal storage is underexplored.
    \item \textbf{Real-time pointing optimization:} Most sun-pointing vs.\ thrust-pointing work is offline trajectory planning; onboard MPC implementations are limited.
\end{enumerate}

\subsection{Comparison: AI Search vs.\ Database Search}

The AI-assisted search using multiple API-accessible databases (arXiv, CrossRef, OpenAlex, CORE, Semantic Scholar) found papers that complement the Web of Science searches already conducted:

\begin{itemize}
    \item \textbf{Overlap:} Both approaches found papers on attitude estimation (EKF variants, TRIAD methods) and power scheduling.
    \item \textbf{AI search strengths:}
    \begin{itemize}
        \item Found recent preprints (arXiv 2025--2026) not yet indexed in WoS
        \item Identified theses and conference papers from SmallSat, IEEE Aerospace
        \item Cross-database deduplication by DOI
        \item Better coverage of open-access repositories (HAL, Zenodo, CORE)
    \end{itemize}
    \item \textbf{Database search strengths:}
    \begin{itemize}
        \item Boolean query precision for complex multi-term searches
        \item Better coverage of older/paywalled journal literature
        \item Citation metrics and impact factors readily available
    \end{itemize}
    \item \textbf{Unique AI findings:} The inertial morphing / Dzhanibekov control literature (Castro-Fern\'{a}ndez, Trivailo, Ostanin) was more effectively discovered through semantic search than Boolean queries.
\end{itemize}

\end{document}
